% ------------------------------------------------------------------
%  Capitolo 18 — Motivazione e scopo
%  Parte IV
%  Ricerca di riferimento: ricerca 01 §10; 03 §2.3, §3.3
%  Voci bibliografiche principali: ryan2000, wigfield2000, hulleman2009, yeager2014, sisk2018, macnamara2023, crede2017, gruber2014
%  Epigrafe: ✔ verificata sul testo (vedi ricerca 03 e registro, ricerca 00)
%  Scheda del capitolo: ricerca/06_indice-ragionato.md, capitolo 18
% ------------------------------------------------------------------
\chapter{Motivazione e scopo}
\label{cap:motivazione}

\epigrafe[ricordano tutto ciò che sta loro a cuore]{omnia, quae curant, meminerunt}{Cicerone, Cato maior de senectute 21}

\iniziale{Q}{uesto capitolo} \segnaposto{apertura: da scrivere — vedi la scheda nell'indice ragionato}.

\section{Si ricorda ciò che sta a cuore}

\segnaposto{da scrivere}

\section{Autonomia, competenza, relazione}

\segnaposto{da scrivere}

\section{Il valore di ciò che si studia}

\segnaposto{da scrivere}

\section{Mentalità e grinta: meno di quanto si dice}

\segnaposto{da scrivere}

\section{La curiosità}

\segnaposto{da scrivere}

\begin{daricordare}
\segnaposto{il principio del capitolo in due righe}
\end{daricordare}

\begin{domande}
  \item \segnaposto{domanda di recupero su questo capitolo}
  \item \segnaposto{domanda di applicazione a una materia che studi}
  \item \segnaposto{domanda di ripresa su un capitolo precedente}
\end{domande}

\begin{sintesi}
  \item \segnaposto{punto di sintesi}
\end{sintesi}

\finecapitolo
